\section{Discussion}\label{sec:discussion}

{\color{red}
The empirical comparison shows why financial fidelity benefits from
a profile of interpretable discrepancies.
AIL has the smallest aggregate deviation, yet its marginal distribution,
long-lag volatility clustering, and aggregational Gaussianity remain
below real--real gap parity at the reported interval level.
MSV has the highest volatility-clustering score, while SBBTS improves
on the controls for aggregational Gaussianity and conditional tails
without recovering their marginal or dependence structure.
A single ranking would conceal these distinctions.
The individual metrics instead identify which properties require
further examination before a generated panel is used in a particular
financial application.
The real--real baseline makes these discrepancies interpretable
relative to the heterogeneity of the observed ticker corpus.
It avoids treating any nonzero distance as evidence of a generator
failure, while retaining information about the magnitude of the
synthetic--real gap.
However, the reference of $0.5$ is an algebraic gap-parity condition,
not a universal threshold for financial adequacy.
An interval containing this value does not establish equivalence,
and a high score is not a substitute for examining the underlying
features.
In the present comparison, AIL's M4 interval includes the reference
but also permits appreciable discrepancies.
Its M3 interval ends close to the reference and is comparatively wide,
so the distinction should not be overstated.
The controls provide complementary checks of the framework.
GBM exposes the response to independent Gaussian returns, while MSV
tests whether a generator calibrated to persistent absolute-return
dependence receives greater agreement on M2.
MSV's performance is therefore informative as a positive control,
but its calibration to the same historical corpus limits any claim
of independent validation.
The empirical comparison is an in-sample fidelity assessment over the
full historical window, including the SBBTS early-stopping dates.
Although those dates are excluded from parameter fitting, they influence
model selection and therefore do not constitute an untouched evaluation
period.
The results describe fidelity to the observed corpus and do not
establish out-of-sample generalisation.
The completed numerical analyses separate two sources of variability.
The outer bootstrap varies the composition of the available ticker
corpora, whereas the 20-seed analysis varies the inner matched draws
on fixed corpora.
The leading generator for each metric and the aggregate ordering
are unchanged across these seeds, but absolute scores still vary.
Neither analysis includes refitting or regenerating the models,
and neither measures uncertainty over new trading dates or market
regimes.
The intervals are pointwise summaries; no simultaneous coverage
claim across the full set of metric--generator comparisons is made.
Several further limitations remain.
The sample covers one equity universe and one historical interval,
including an unusually turbulent market episode.
The selected feature maps retain specific distributional and temporal
signatures, so agreement on them cannot establish agreement of the
full joint return distribution.
Some metric parameters were chosen during development with AIL
outputs visible. The 19-specification analysis retains the complete M1,
M4 and aggregate orderings across five paired seeds, while the scores
still respond to the chosen parameters. This supports stability of the
observed rankings within the examined grid; it does not remove that
selection history. The examined ranges and three
selected interactions do not cover every structural metric choice.
Finally, stylized-fact fidelity does not by itself demonstrate
downstream trading utility, privacy, or suitability for risk estimation.
These limitations define the role of the framework as a diagnostic
layer in a broader evaluation.
Independent market periods, repeated generator fits, broader structural
sensitivity analyses, and application-specific utility tests would provide
complementary evidence.
Within the present scope, the contribution is a transparent comparison
of financially meaningful discrepancies against an observed real-data
reference, with corpus uncertainty and numerical stability reported
separately.
}
